% !PRINTED 677
\weberpage{677}
%<<CONT
in which $\mathfrak{U}$ and $\mathfrak{U}'$ are equivalent $n$-gons without common divisor. An integral function $\omega$ of $z$ is characterised by the fact that it becomes infinite at no point which is not contained in $\mathfrak{U}$. Consequently such a function $\omega$ can be represented in the following way:
\weq{2}{\omega = \dfrac{\mathfrak{A}}{\mathfrak{U}^r},}
in which $r$ is a positive exponent. In (2) the numerator and the denominator may have common divisors. We may, however, assume that $\mathfrak{A}$ is not divisible by $\mathfrak{U}$. Then in (2) the exponent $r$ has the least possible value. \gesperrt{This number $r$ shall be called the exponent of $\omega$ in $z$.}

1. \gesperrt{If we set $z' = 1:z$, then by (1), (2):}
\[
z' = \frac{\mathfrak{U}}{\mathfrak{U}'}, \quad \omega' = \omega z'^r = \frac{\mathfrak{A}}{\mathfrak{U}'^r},
\]
and if $\mathfrak{A}$ is not divisible by $\mathfrak{U}'$, then $r$ is the exponent of $\omega'$ with respect to $z'$. The assumption that $\mathfrak{A}$ is not divisible by $\mathfrak{U}'$ means that $\omega$ is not divisible by $z$, and from this it follows that $\omega'$ too is not divisible by $z'$.

2. The exponent of an integral rational function of $z$ of degree $m$
\[
x = a_0 + a_1 z + \cdots + a_m z^m
\]
is equal to $m$. For this function becomes infinite of the $m$th order at every point contained in $\mathfrak{U}$. The exponent of the product $x\omega$ is equal to $r + m$.

3. The exponent of a sum
\[
\omega = c_1\omega_1 + c_2\omega_2 + \cdots + c_s\omega_s
\]
is at most equal to the greatest of the exponents of $\omega_1, \omega_2, \ldots, \omega_s$.

4. We determine a sequence of integral functions of $z$
\[
\lambda_1, \lambda_2, \lambda_3, \ldots
\]
by the following recurrent determination:
\begin{itemize}
\item[1)] $\lambda_1$ constant (exponent $r_1 = 0$),
\item[2)] $\lambda_2$ a non-rational integral function of $z$ with the smallest possible exponent $r_2$,
\item[3)] $\lambda_s$ an integral function of $z$ not contained in the form
\weq{3}{x_1\lambda_1 + x_2\lambda_2 + \cdots + x_{s-1}\lambda_{s-1}}
with the smallest possible exponent $r_s$,
\end{itemize}
%CONT>>
